
%allgemeine Formatangaben
\documentclass[
 a4paper, 										% Papierformat
 12pt,												% Schriftgr��e
 ngerman, 										% f�r Umlaute, Silbentrennung etc.
 titlepage,										% es wird eine Titelseite verwendet
 bibliography=totoc,					% Literaturverzeichnis im Inhaltsverzeichnis auff�hren
 listof=totoc,								% Verzeichnisse im Inhaltsverzeichnis auff�hren
 oneside, 										% einseitiges Dokument
 captions=nooneline,					% einzeilige Gleitobjekttitel ohne Sonderbehandlung wie mehrzeilige Gleitobjekttitel behandeln
 numbers=noenddot,						% �berschriften-??Nummerierung ohne Punkt am Ende
 parskip=half									% zwischen Abs�tzen wird eine halbe Zeile eingef�gt
 ]{scrbook}


\newcommand{\titel}{Entwicklung eines Bilderkennungssystems zum finden von Duplikaten in einer Bilddatenbank bei der Registrierung neuer Bilder}
\newcommand{\abschlussart}{Bachelor of Science (B.Sc.)}
\newcommand{\arbeit}{Bachelorarbeit}
\newcommand{\hochschule}{Hochschule f�r Technik, Wirtschaft und Kultur Leipzig}
\newcommand{\fakultaet}{Fakult�t Informatik und Medien}
\newcommand{\autor}{Martino Thomann}
\newcommand{\studiengang}{Bachelorstudiengang Medieninformatik}
\newcommand{\matrikelnr}{75484}
\newcommand{\erstgutachter}{Prof. Dr. Sibylle Schwarz}
\newcommand{\zweitgutachter}{B.Sc. Marc Bellmann}
\newcommand{\ort}{Leipzig}
 			% einbinden von pers�nlichen Daten



% Anpassung an Landessprache
\usepackage[ngerman]{babel}
\usepackage{natbib}
 
% Verwenden von Sonderzeichen und Silbentrennung
\usepackage[latin1]{inputenc}	
\usepackage[T1]{fontenc}			
\usepackage{textcomp} 																% Euro-Zeichen und andere
\usepackage[babel,german=quotes]{csquotes}						% Anf�hrungszeichen
\RequirePackage[ngerman=ngerman-x-latest]{hyphsubst} 	% erweiterte Silbentrennung

% Befehle aus AMSTeX f�r mathematische Symbole z.B. \boldsymbol \mathbb
\usepackage{amsmath,amsfonts}

% Zeilenabst�nde und Seitenr�nder 
\usepackage{setspace}
\usepackage{geometry}

% Einbinden von JPG-Grafiken
\usepackage{graphicx}

% zum Umflie�en von Bildern
% Verwendung unter http://de.wikibooks.org/wiki/LaTeX-Kompendium:_Baukastensystem#textumflossene_Bilder
\usepackage{floatflt}

% Verwendung von vordefinierten Farbnamen zur Colorierung
% Palette und Verwendung unter http://kitt.cl.uzh.ch/kitt/CLinZ.CH/src/Kurse/archiv/LaTeX-Kurs-Farben.pdf
\usepackage[usenames,dvipsnames]{color} 

% Tabellen
\usepackage{array}
\usepackage{longtable}

% einfache Grafiken im Code
% Einf�hrung unter http://www.math.uni-rostock.de/~dittmer/bsp/pstricks-bsp.pdf
\usepackage{pstricks}

% Quellcodeansichten
\usepackage{verbatim}
\usepackage{moreverb} 											% f�r erweiterte Optionen der verbatim Umgebung
% Befehle und Beispiele unter http://www.ctex.org/documents/packages/verbatim/moreverb.pdf
\usepackage{listings} 											% f�r angepasste Quellcodeansichten siehe
% Kurzeinf�hrung unter http://blog.robert-kummer.de/2006/04/latex-quellcode-listing.html

% Glossar und Abbildungsverzeichnis
\usepackage[
nonumberlist, %keine Seitenzahlen anzeigen
acronym,      %ein Abk�rzungsverzeichnis erstellen
toc          %Eintr�ge im Inhaltsverzeichnis
]      %im Inhaltsverzeichnis auf section-Ebene erscheinen
{glossaries}

% verlinktes und Farblich angepasstes Inhaltsverzeichnis
\usepackage[pdftex,
colorlinks=true,
linkcolor=InterneLinkfarbe,
urlcolor=ExterneLinkfarbe]{hyperref}
\usepackage[all]{hypcap}

% URL verlinken, lange URLs umbrechen
\usepackage{url}

% sorgt daf�r, dass Leerzeichen hinter parameterlosen Makros nicht als Makroendezeichen interpretiert werden
\usepackage{xspace}

% Beschriftungen f�r Abbildungen und Tabellen
\usepackage{caption}

% Entwicklerwarnmeldungen entfernen
\usepackage{scrhack}
					% einbinden der verwendeten Latex-Pakete


\onehalfspacing 							% 1,5facher Zeilenabstand

\definecolor{InterneLinkfarbe}{rgb}{0.1,0.1,0.3} 	% Farbliche Absetzung von externen Links
\definecolor{ExterneLinkfarbe}{rgb}{0.1,0.1,0.7}	% Farbliche Absetzung von internen Links

% Einstellungen f�r Fu�noten:
\captionsetup{font=footnotesize,labelfont=sc,singlelinecheck=true,margin={5mm,5mm}}

% Stil der Quellenangabe
\bibliographystyle{alphadin}

%Ausschluss von Schusterjungen
\clubpenalty = 10000
%Ausschluss von Hurenkindern
\widowpenalty = 10000

% Befehle, die Umlaute ausgeben, f�hren zu Fehlern, wenn sie hyperref als Optionen �bergeben werden
\hypersetup{
%    pdftitle={\titel \untertitel},
%    pdfauthor={\autor},
%    pdfcreator={\autor},
%    pdfsubject={\titel \untertitel},
%    pdfkeywords={\titel \untertitel},
}

% Beispiel f�r eine Listings-Codeumbebungen
% Bei mehreren Definitionen empfielt sich das auslagern in eine externe Datei
\lstloadlanguages{Java,HTML}
\lstset{
	frame=tb,
	framesep=5pt,
	basicstyle=\footnotesize\ttfamily,
	showstringspaces=false,
	keywordstyle=\ttfamily\bfseries\color{CadetBlue},
	identifierstyle=\ttfamily,
	stringstyle=\ttfamily\color{OliveGreen},
	commentstyle=\color{GrayBlue},
	rulecolor=\color{Gray},
	xleftmargin=5pt,
	xrightmargin=5pt,
	aboveskip=\bigskipamount,
	belowskip=\bigskipamount
} 

%Den Punkt am Ende jeder Beschreibung deaktivieren
\renewcommand*{\glspostdescription}{}

% Empfehlung: Abkuerzungsverzeichnis und Glossar sind in Graduierunsarbeiten
% nicht zwingend notwendig

% %Glossar-Befehle anschalten
% \makeglossaries
% \glsenablehyper
% \input{Header/Abkuerzungen}
% \input{Header/Glossar}


\areaset[0.5cm]{16cm}{24cm}


\begin{document}

%\include{Deckblatt/titel}						% Deckblatt der vorliegenden Arbeit
% Die Daten werden in Header/Metadaten.tex eingstellt
\begin{titlepage}
\begin{Large}
\begin{center}

\includegraphics[width=0.5\textwidth]{Abbildungen/HTWK_Zusatz_de_V_Black_sRGB.png}
% \textbf{\hochschule\ \ort}\\[5pt]
% \fachbereich\\
% \studiengang\\
% \vskip 1cm
% \arbeit\\
% zur Erlangung der akademischen Grades
\\[8pt]

\arbeit\\[-1mm]
zur Erlangung des akademischen Grades\\[3mm]
\abschlussart\\[3mm]
im \studiengang\\[-1mm]
der \fakultaet%\\[-1mm]
%der \hochschule  %% kann entfallen
\\
\vfill
{\LARGE\bfseries \titel \par}
\vfill

vorgelegt von\\
 \autor%\\
% Matrikelnummer: \matrikelnr  
\\[8pt]
\ort, den \today

\end{center}
\vfill
\begin{tabular}{ll}
Erstpr�fer: & \erstgutachter\\
Zweitpr�fer: & \zweitgutachter\\
\end{tabular}
\end{Large}
\end{titlepage}

\frontmatter												% Seitenzählerstart vor dem Text


\chapter*{Zusammenfassung}
\label{sec:Zusammenfassung}
In dieser Arbeit wird der Entwurf eines Bilderkennungssystems beschrieben. 
Das System ist in der Lage  zu erkennen, ob es sich bei einem neuen Bild um Duplikat handelt, das bereits in einer Datenbank mit registrierten Bildern vorhanden ist.
Um Duplikate zu erkennen, deren Motiv im Vergleich zum Original modifiziert, aber immer noch als Kopie erkennbar sind, wird ein feature-based Bilderkennungs-Algorithmus eingesetzt. 
Innerhalb der wurden SIFT, ORB und BRIEF miteinander im Anwendungsfall von T-Shirt Designs bei Spreadshirt innerhalb von verschiedenen Szenarien getestet und miteinander verglichen.
						% Abstrakt

%% Empfehlung: Auf Danksagungen + Vorwort in wissenschaftlichen Arbeiten eher verzichten -- ggf. die entsprechenden Dateien im Ordner Inhalt anlegen und hier auskommentieren (PKW)
%% \include{Inhalt/Danksagung}					% Danksagung
%% \include{Inhalt/Vorwort}						% Vorwort

\tableofcontents										% Inhaltsverzeichnis

\mainmatter													% Seitenzählerstart Haupttext


\chapter{Einleitung und Intension}
\label{sec:EinleitungUndIntension}
Was m�chte ich sagen? Adawdadwadasda

\chapter{Erste �berschrift}
\label{sec:ErsteUeberschrift}
Hier steht der erste Text.

\section{Erstes Unterkapitel}
\label{sec:ErstesUnterkapitel}
Und hier wird es schon spezifischer.


\section{Zweites Unterkapitel}
\label{sec:ZweitesUnterkapitel}

In der Regel reichen zwei Gliederungsebenen f�r Graduierungsarbeiten aus.					% hier steht der eigentliche Text der Arbeit

\pagenumbering{Roman}
\appendix	
												
\bibliography{Literatur/Literatur} 	% Literaturverzeichnis

% Empfehlung: Verzeichnisse an das Ende der Arbeit (PKW)
% Ein Glossar ist eher untypisch fuer Graduierungsarbeiten
% Abkuerzungsverzeichnis an das Ende der Arbeit; Abkuerzungen und Begriffe muessen in jedem
% Fall im Haupttext der Arbeit beim ersten Vorkommen definiert werden

%% \printglossary[title=Glossar] 			% Glossar Einträge in Header/Glossar.tex vornehmen
%%\printglossary[type=\acronymtype,title=Abk\"urzungsverzeichnis]	% Abkürzungsverzeichnis Einträge in Header/Abkuerzungen vornehmen

\listoffigures											% Abbildungsverzeichnis
\listoftables												% Tabellenverzeichnis


\pagestyle{empty}

\chapter*{Selbst�ndigkeitserkl�rung}
\label{sec:Selbst�ndigkeitserkl�rung}

Ich erkl�re hiermit, dass ich die vorliegende Graduierungsarbeit ohne Hilfe Dritter und ohne Benutzung anderer als der angegebenen Quellen und Hilfsmittel verfasst habe. Alle den benutzten Quellen w�rtlich oder sinngem�� entnommene Stellen sind als solche einzeln kenntlich gemacht.

Diese Arbeit ist bislang keiner anderen Pr�fungsbeh�rde vorgelegt und auch nicht ver�ffentlicht worden.

Ich bin mir bewusst, dass eine falsche Erkl�rung rechtliche Folgen haben wird.\\[1.5cm]


\begin{minipage}{0.7\textwidth}
\ort, \today \hfill  Unterschrift
\end{minipage}





\end{document}